\documentclass{article}
\usepackage[T1]{fontenc}
\usepackage{lmodern}
\usepackage{graphicx}
\usepackage{amsmath,amssymb}
\usepackage[a4paper,margin=2cm]{geometry}
\usepackage[absolute,overlay]{textpos}
\setlength{\TPHorizModule}{1mm}
\setlength{\TPVertModule}{1mm}
\usepackage[indonesian]{babel}
\usepackage{hyperref}
\setlength{\emergencystretch}{2em}
% unit: Halaman 6

\begin{document}

% === Halaman 6 (llm) ===

4. $\langle Z_n, \oplus \rangle$ : grup dengan himpunan \textit{integer} modulo $n$, yaitu $Z_n = \{0, 1, 2, ..., n - 1\}$ dan $\oplus$ adalah operasi penjumlahan modulo $n$.

Contoh: $n = 5$, $Z_n = \{0, 1, 2, 3, 4\}$, $(3 \oplus 4) = (3 + 4) \pmod 5 = 2$

$\langle Z_p, \oplus \rangle$ : grup dengan himpunan \textit{integer} modulo $p$, $p$ adalah bilangan prima, yaitu $Z_p = \{0, 1, 2, ..., p - 1\}$ dan $\oplus$ adalah operasi penjumlahan modulo $p$.

$\langle Z_p^*, \otimes \rangle$ : dengan himpunan integer bukan nol, $p$ adalah bilangan prima, yaitu $Z_p^* = \{1, 2, ..., p - 1\}$ dan $\otimes$ adalah operasi perkalian modulo $p$.

\end{document}
